\documentclass[12pt]{article}
\usepackage{fullpage}
\usepackage[T1]{fontenc}
\usepackage[utf8]{inputenc}
\usepackage{lmodern}
\usepackage{microtype}
\usepackage{amsmath,amssymb,amsthm}
\usepackage{hyperref}

\newtheorem{theorem}{Theorem}
\newtheorem{proposition}{Proposition}
\newtheorem{lemma}{Lemma}
\newtheorem{corollary}{Corollary}
\theoremstyle{definition}
\newtheorem{definition}{Definition}
\newtheorem{remark}{Remark}

\DeclareMathOperator{\inv}{inv}
\DeclareMathOperator{\TV}{TV}
\DeclareMathOperator{\Var}{Var}
\newcommand{\TW}{\mathrm{TW}_2}
\newcommand{\E}{\mathbb{E}}
\newcommand{\PP}{\mathbb{P}}
\newcommand{\R}{\mathbb{R}}
\newcommand{\eps}{\varepsilon}

\begin{document}

\title{KPZ universality for longest increasing subsequences \\ of Mallows--distributed permutations}
\author{}
\date{}
\maketitle

\begin{abstract}
We study the longest increasing subsequence $L(\pi)$ when $\pi$ is drawn from the
Mallows measure $\mu_{n,q}$ on $S_n$, $\mu_{n,q}(\pi)\propto q^{\inv(\pi)}$. Writing
$\eps_n=1-q_n$, we prove three complementary statements that together characterize the
fluctuation behaviour.
\begin{enumerate}
\item[(A)] (\emph{KPZ regime}) If $q_n\to1$ with $n^{3/2}\eps_n\to0$, then
$\big(L(\pi)-2\sqrt n\big)/n^{1/6}\xrightarrow{d}F_{\TW}$, the GUE Tracy--Widom law,
with the Baik--Deift--Johansson centering and scaling.
\item[(B)] (\emph{Sharpness}) $\TV(\mu_{n,q_n},U_n)\to0$ if and only if
$n^{3/2}\eps_n\to0$; and $\TV(\mu_{n,q_n},U_n)\to1$ when $n^{3/2}\eps_n\to\infty$.
\item[(C)] (\emph{Ballistic crossover}) For fixed $q\in(0,1)$, there is $c(q)>0$ with
$\PP_{\mu_{n,q}}(L(\pi)\ge c(q)\,n)\to1$; hence $L\asymp n$ and both the
$\sqrt n$ scale and the Tracy--Widom law are destroyed.
\end{enumerate}
All bounds are explicit and non-asymptotic. The only nontrivial external input is the
Baik--Deift--Johansson theorem for the uniform measure.
\end{abstract}

\section{Setup and the external input}\label{sec:setup}

\subsection{Definitions}
For $\pi\in S_n$, $\inv(\pi)=\#\{(i,j):i<j,\ \pi(i)>\pi(j)\}$, $L(\pi)$ is the length of
a longest increasing subsequence, and $D(\pi)$ the length of a longest \emph{decreasing}
subsequence. For $q\in(0,1]$ the \emph{Mallows measure} on $S_n$ is
\begin{equation}\label{eq:mallows}
\mu_{n,q}(\pi)=\frac{q^{\inv(\pi)}}{Z_n(q)},\qquad
Z_n(q)=\sum_{\sigma\in S_n}q^{\inv(\sigma)}=\prod_{i=1}^{n}\frac{1-q^{i}}{1-q},
\end{equation}
the last equality being the classical identity $Z_n(q)=[n]_q!$ \cite[\S1.3]{Stanley}.
We write $U_n=\mu_{n,1}$ (uniform), $\eps=1-q$, and $\PP_\mu,\E_\mu$ for probability and
expectation under a measure $\mu$.

\subsection{The Lehmer code and its exact law}\label{sec:code}
For $\pi\in S_n$ define $c(\pi)=(c_1,\dots,c_n)$ by
$c_i(\pi)=\#\{\,j>i:\pi(j)<\pi(i)\,\}$. The map $\pi\mapsto c(\pi)$ is a bijection of
$S_n$ onto $\prod_{i=1}^n\{0,1,\dots,n-i\}$ (its inverse processes $i=1,\dots,n$ and
sets $\pi(i)$ to be the $(c_i{+}1)$-st smallest unused value; at step $i$ exactly
$n-i+1$ values remain, so $c_i\in\{0,\dots,n-i\}$). Moreover, since $c_i$ counts exactly
the inversions whose smaller index is $i$ and every inversion is counted once,
\begin{equation}\label{eq:invsum}
\inv(\pi)=\sum_{i=1}^{n}c_i(\pi).
\end{equation}

\begin{lemma}[Product form of the Mallows code]\label{lem:product}
Under $\mu_{n,q}$ the coordinates $c_1,\dots,c_n$ are independent, and $c_i$ has law
\begin{equation}\label{eq:nu}
\nu^{(q)}_{m}(k)=\frac{q^{k}}{1+q+\cdots+q^{m-1}},\qquad k\in\{0,1,\dots,m-1\},
\end{equation}
on the support of size $m=n-i+1$. Under $U_n$ the $c_i$ are independent and $c_i$ is
uniform on $\{0,\dots,n-i\}$, i.e.\ $\nu^{(1)}_m$ is uniform on $m$ points.
\end{lemma}

\begin{proof}
By \eqref{eq:mallows} and \eqref{eq:invsum},
$\mu_{n,q}(\pi)=Z_n(q)^{-1}\prod_{i=1}^n q^{c_i(\pi)}$. Since $\pi\mapsto c(\pi)$ is a
bijection onto $\prod_i\{0,\dots,n-i\}$, the pushforward of $\mu_{n,q}$ under $c$ is the
product measure whose $i$-th factor gives mass $\propto q^k$ to $k\in\{0,\dots,n-i\}$;
normalizing gives \eqref{eq:nu} with $m=n-i+1$. Consistency with \eqref{eq:mallows}:
$\prod_{i=1}^n(1+q+\cdots+q^{n-i})=\prod_{m=1}^n[m]_q=Z_n(q)$. For $q=1$ each factor is
uniform.
\end{proof}

\subsection{The Baik--Deift--Johansson theorem}\label{sec:bdj}
We use exactly one external result.

\begin{theorem}[Baik--Deift--Johansson \cite{BDJ}]\label{thm:BDJ}
If $\pi$ is uniform on $S_n$, then for every $s\in\R$,
\[
\PP_{U_n}\!\big(L(\pi)\le 2\sqrt n+s\,n^{1/6}\big)\xrightarrow[n\to\infty]{}F_{\TW}(s),
\]
where $F_{\TW}$ is the GUE Tracy--Widom distribution function, which is continuous on
$\R$.
\end{theorem}

Continuity (indeed real-analyticity) of $F_{\TW}$ is classical \cite{TW}; it is the only
property of $F_{\TW}$ we use.

\section{Theorem A: the KPZ regime $n^{3/2}\eps_n\to0$}\label{sec:A}

\subsection{A coordinatewise $\chi^2$ bound}
For probability measures $\mu,\nu$ on a finite set with $\nu>0$,
$\chi^2(\mu\,\|\,\nu)=\sum_x \mu(x)^2/\nu(x)-1$ and
$\TV(\mu,\nu)=\tfrac12\sum_x|\mu(x)-\nu(x)|$.

\begin{lemma}\label{lem:chi2}
For every integer $m\ge1$ and $q\in(0,1]$, with $\eps=1-q$,
\begin{equation}\label{eq:chi2bound}
\chi^2\!\big(\nu^{(q)}_m\,\big\|\,\nu^{(1)}_m\big)\;\le\;\frac{\eps^2\,m^2}{12\,q^{2(m-1)}}.
\end{equation}
\end{lemma}

\begin{proof}
For $q=1$ both sides vanish; assume $q\in(0,1)$ and put $S=\sum_{k=0}^{m-1}q^{k}$. Since
$\nu^{(1)}_m(k)=1/m$,
\[
\chi^2\!\big(\nu^{(q)}_m\,\big\|\,\nu^{(1)}_m\big)
=\frac{m\sum_{k=0}^{m-1}q^{2k}}{S^2}-1
=\frac{m\sum_{k=0}^{m-1}q^{2k}-\big(\sum_{k=0}^{m-1}q^{k}\big)^2}{S^2}.
\]
By the Lagrange identity $m\sum_k a_k^2-(\sum_k a_k)^2=\sum_{0\le i<j\le m-1}(a_i-a_j)^2$
with $a_k=q^k$, the numerator equals $\sum_{0\le i<j\le m-1}(q^{i}-q^{j})^2$. For $i<j$,
$q^{i}-q^{j}=q^{i}(1-q^{j-i})$ with $0\le q^{i}\le1$ and, by
$1-q^{d}=(1-q)\sum_{t=0}^{d-1}q^{t}\le d(1-q)=d\eps$, also $0\le 1-q^{j-i}\le(j-i)\eps$;
hence $(q^i-q^j)^2\le (j-i)^2\eps^2$ and
\[
\sum_{0\le i<j\le m-1}(q^{i}-q^{j})^2\le \eps^2\sum_{d=1}^{m-1}(m-d)d^2
=\eps^2\,\frac{m^2(m^2-1)}{12}\le \eps^2\,\frac{m^4}{12}.
\]
For the denominator, each term of $S$ is $\ge q^{m-1}$, so $S\ge mq^{m-1}$ and
$S^2\ge m^2q^{2(m-1)}$. Dividing gives \eqref{eq:chi2bound}.
\end{proof}

\subsection{Tensorization and total variation}
\begin{proposition}\label{prop:TVupper}
For all $n\ge1$ and $q\in(0,1]$ with $\eps=1-q$,
\begin{equation}\label{eq:TVupper}
\TV\big(\mu_{n,q},U_n\big)\le
\tfrac12\sqrt{\,\exp\!\Big(\tfrac{\eps^2}{12}\sum_{m=1}^{n}\tfrac{m^2}{q^{2(m-1)}}\Big)-1\,}.
\end{equation}
If moreover $2n\eps\le1$, then
\begin{equation}\label{eq:TVclean}
\TV\big(\mu_{n,q},U_n\big)\le \tfrac12\sqrt{\exp\!\Big(\tfrac{n^3\eps^2}{6\,(1-2n\eps)}\Big)-1}.
\end{equation}
Consequently $\TV(\mu_{n,q_n},U_n)\to0$ whenever $n^{3/2}\eps_n\to0$.
\end{proposition}

\begin{proof}
By Lemma~\ref{lem:product} the laws of $c=(c_1,\dots,c_n)$ under $\mu_{n,q}$ and $U_n$
are product measures $\bigotimes_i\nu^{(q)}_{n-i+1}$ and $\bigotimes_i\nu^{(1)}_{n-i+1}$,
and $\pi\mapsto c(\pi)$ is a bijection, so $\TV(\mu_{n,q},U_n)$ equals the total
variation of these products. The $\chi^2$ divergence tensorizes:
$1+\chi^2(\bigotimes_i\alpha_i\,\|\,\bigotimes_i\beta_i)=\prod_i(1+\chi^2(\alpha_i\,\|\,\beta_i))$,
since $\sum_x\prod_i\alpha_i(x_i)^2/\prod_i\beta_i(x_i)=\prod_i\sum_{x_i}\alpha_i(x_i)^2/\beta_i(x_i)$.
With $1+t\le e^t$ and Lemma~\ref{lem:chi2} (the index $m=n-i+1$ ranges over
$1,\dots,n$),
\[
1+\chi^2(\mu_{n,q}\,\|\,U_n)=\prod_{m=1}^n\big(1+\chi^2(\nu^{(q)}_m\,\|\,\nu^{(1)}_m)\big)
\le\exp\!\Big(\tfrac{\eps^2}{12}\sum_{m=1}^n\tfrac{m^2}{q^{2(m-1)}}\Big).
\]
The Cauchy--Schwarz bound
$2\TV=\sum_x|\mu-\nu|=\sum_x\frac{|\mu-\nu|}{\sqrt\nu}\sqrt\nu\le\sqrt{\sum_x\frac{(\mu-\nu)^2}{\nu}}
=\sqrt{\chi^2(\mu\,\|\,\nu)}$ gives \eqref{eq:TVupper}.

If $2n\eps\le1$ then $q^{2(m-1)}\ge q^{2(n-1)}=(1-\eps)^{2(n-1)}\ge 1-2(n-1)\eps\ge1-2n\eps>0$
(Bernoulli), and $\sum_{m=1}^n m^2\le n^3$, so the exponent is at most
$\frac{\eps^2 n^3}{12(1-2n\eps)}\le \frac{n^3\eps^2}{6(1-2n\eps)}$, which is
\eqref{eq:TVclean}. If $n^{3/2}\eps_n\to0$, then $n\eps_n\le n^{3/2}\eps_n\to0$ so
eventually $2n\eps_n\le1$ and $1-2n\eps_n\to1$, while $n^3\eps_n^2=(n^{3/2}\eps_n)^2\to0$;
hence the right-hand side of \eqref{eq:TVclean} $\to0$.
\end{proof}

\subsection{Conclusion of Theorem A}
\begin{theorem}\label{thm:A}
If $q_n\in(0,1]$, $q_n\to1$ and $n^{3/2}(1-q_n)\to0$, and $\pi\sim\mu_{n,q_n}$, then
$\dfrac{L(\pi)-2\sqrt n}{n^{1/6}}\xrightarrow{d}F_{\TW}$.
\end{theorem}

\begin{proof}
Fix $s\in\R$ and $A_n=\{\pi:L(\pi)\le 2\sqrt n+s\,n^{1/6}\}\subseteq S_n$. By the
variational definition of total variation,
$|\PP_{\mu_{n,q_n}}(A_n)-\PP_{U_n}(A_n)|\le \TV(\mu_{n,q_n},U_n)\to0$ by
Proposition~\ref{prop:TVupper}. By Theorem~\ref{thm:BDJ}, $\PP_{U_n}(A_n)\to F_{\TW}(s)$.
Hence $\PP_{\mu_{n,q_n}}(A_n)\to F_{\TW}(s)$ for every $s$. As $F_{\TW}$ is continuous,
pointwise convergence of distribution functions at all $s$ is convergence in
distribution.
\end{proof}

\begin{remark}
For $q_n>1$ ($\eps_n<0$) the same proof applies with $q^{2(m-1)}$ in
Lemma~\ref{lem:chi2} replaced by $1$ (then $q^i\ge1$ and $S\ge m$), so the conclusion
holds throughout $|1-q_n|=o(n^{-3/2})$.
\end{remark}

\section{Theorem B: sharpness of the threshold}\label{sec:B}

We show $n^{3/2}\eps_n\to0$ is exactly the coupling threshold by separating the two
measures with the statistic $X=\inv$. We first record two facts about
$\nu^{(t)}_m$.

\begin{lemma}\label{lem:moments}
Let $c'\sim\nu^{(t)}_m$, $m\ge2$. Then:
\begin{enumerate}
\item[(i)] $\nu^{(1)}_m$ has mean $\tfrac{m-1}{2}$ and variance $\tfrac{m^2-1}{12}$.
\item[(ii)] $g_m(t):=\E_{\nu^{(t)}_m}[c']$ is nondecreasing in $t$ on $(0,1]$
(equivalently, $g_m(1)-g_m(t)\ge0$ is nondecreasing in $\eps=1-t$), and
\[
g_m(1)-g_m(t)=\int_t^1 \frac{\Var_{\nu^{(s)}_m}(c')}{s}\,ds\qquad(0<t\le1).
\]
\item[(iii)] For $s\in[\,1-\tfrac1{2m},1\,]$ one has $\Var_{\nu^{(s)}_m}(c')\ge\tfrac{m^2-1}{48}$.
\end{enumerate}
\end{lemma}

\begin{proof}
(i) Standard moments of the uniform law on $\{0,\dots,m-1\}$.

(ii) Write $\theta=\log t\le0$ and $h(\theta)=\log\sum_{k=0}^{m-1}e^{\theta k}$, the
cumulant generating function of $\nu^{(1)}_m$. Then $\nu^{(t)}_m(k)=e^{\theta k-h(\theta)}$
is the exponential tilt of $\nu^{(1)}_m$ by $\theta$, so
$\E_{\nu^{(t)}_m}[c']=h'(\theta)=:G(\theta)$ and
$\Var_{\nu^{(t)}_m}(c')=h''(\theta)\ge0$. Hence $G$ is nondecreasing in $\theta$, and
since $\theta=\log t$ is increasing in $t$, $g_m(t)=G(\log t)$ is nondecreasing in $t$
on $(0,1]$; equivalently $g_m(1)-g_m(t)\ge0$ and is nondecreasing in $\eps=1-t$. The
chain rule gives the integral formula, since
\[
g_m'(s)=\frac{d}{ds}G(\log s)=\frac{h''(\log s)}{s}=\frac{\Var_{\nu^{(s)}_m}(c')}{s},
\qquad\text{so}\qquad
g_m(1)-g_m(t)=\int_t^1 g_m'(s)\,ds=\int_t^1\frac{\Var_{\nu^{(s)}_m}(c')}{s}\,ds.
\]

(iii) Use $\Var_{\nu^{(s)}_m}(c')=\tfrac12\sum_{k,l=0}^{m-1}(k-l)^2\,\nu^{(s)}_m(k)\,\nu^{(s)}_m(l)$,
the standard pairwise representation of the variance. For $s\in[1-\tfrac1{2m},1]$ we
have $s^{m-1}\ge 1-(m-1)(1-s)\ge 1-\tfrac{m-1}{2m}\ge\tfrac12$ (Bernoulli), and with
$S_s=\sum_{k=0}^{m-1}s^k\le m$,
\[
\nu^{(s)}_m(k)=\frac{s^k}{S_s}\ge\frac{s^{m-1}}{m}\ge\frac{1}{2m}\qquad(0\le k\le m-1).
\]
Since $\sum_{k,l=0}^{m-1}(k-l)^2=\tfrac{m^2(m^2-1)}{6}$,
\[
\Var_{\nu^{(s)}_m}(c')\ge\frac12\cdot\frac{1}{(2m)^2}\cdot\frac{m^2(m^2-1)}{6}
=\frac{m^2-1}{48}.\qedhere
\]
\end{proof}

\begin{proposition}\label{prop:TVlower}
There is an absolute constant $C_1=24576$ such that the following hold for all $n\ge4$
and $q\in(0,1]$, $\eps=1-q$.
\begin{enumerate}
\item[(a)] If $\eps\le\tfrac1{2n}$, then
$\displaystyle \TV\big(\mu_{n,q},U_n\big)\ \ge\ 1-\frac{C_1}{(n^{3/2}\eps)^2}$.
\item[(b)] If $\eps>\tfrac1{2n}$, then
$\displaystyle \TV\big(\mu_{n,q},U_n\big)\ \ge\ 1-\frac{4C_1}{n}$.
\end{enumerate}
In particular, if $n^{3/2}\eps_n\to\infty$, then $\TV(\mu_{n,q_n},U_n)\to1$.
\end{proposition}

\begin{proof}
By Lemma~\ref{lem:product}, $X=\inv=\sum_{m=1}^n c'_m$ with $c'_m$ independent,
$c'_m\sim\nu^{(q)}_m$ (relabelling $m=n-i+1$). By Lemma~\ref{lem:moments}(i),
$\E_U[X]=\sum_{m=1}^n\tfrac{m-1}{2}=\tfrac{n(n-1)}{4}$ and
$\sigma_U^2:=\Var_U(X)=\sum_{m=1}^n\tfrac{m^2-1}{12}\le\tfrac{n^3}{12}$.

\emph{Variance control.} Each $c'_m\in\{0,\dots,m-1\}$, an interval of length $m-1$,
so $\Var_\mu(c'_m)\le\tfrac{(m-1)^2}{4}$ (Popoviciu's inequality); hence
$\sigma_\mu^2:=\Var_\mu(X)\le\sum_{m=1}^n\tfrac{(m-1)^2}{4}\le\tfrac{n^3}{12}$, and the
same bound $\sigma_U^2\le\tfrac{n^3}{12}$ holds (it was already computed). Thus
$\sigma_\mu^2+\sigma_U^2\le\tfrac{n^3}{6}$ for every $q\in(0,1]$.

\emph{Mean shift, monotone in $\eps$.} Write $\Delta(\eps):=\E_U[X]-\E_{\mu_{n,1-\eps}}[X]
=\sum_{m=1}^n\big(g_m(1)-g_m(1-\eps)\big)$. By Lemma~\ref{lem:moments}(ii), each summand
is $\ge0$ and nondecreasing in $\eps$, so $\Delta(\eps)$ is nondecreasing in $\eps$.
For $\eps\le\tfrac1{2n}$ we have $1-\eps\ge1-\tfrac1{2m}$ for all $m\le n$, so by
Lemma~\ref{lem:moments}(ii)-(iii), for $m\ge2$,
\[
g_m(1)-g_m(1-\eps)=\int_{1-\eps}^1\frac{\Var_{\nu^{(s)}_m}(c')}{s}\,ds
\ \ge\ \int_{1-\eps}^1\Var_{\nu^{(s)}_m}(c')\,ds\ \ge\ \eps\cdot\frac{m^2-1}{48},
\]
using $1/s\ge1$ for $s\le1$ and that $[1-\eps,1]\subseteq[1-\tfrac1{2m},1]$. Summing
($m=1$ contributes $0$) and using
$\sum_{m=2}^n(m^2-1)=\tfrac{n(n-1)(2n+5)}{6}-(n-1)\ge\tfrac{n^3}{4}$ for $n\ge4$
(elementary),
\begin{equation}\label{eq:Deltalb}
\Delta(\eps)\ \ge\ \frac{\eps}{48}\cdot\frac{n^3}{4}=\frac{\eps\,n^3}{192}
\qquad(0<\eps\le\tfrac1{2n}).
\end{equation}

\emph{Separation.} Fix $q=1-\eps$ and write $\Delta=\Delta(\eps)$. Let
$t^\star=\E_U[X]-\Delta/2$ and $B=\{X\le t^\star\}$. By Chebyshev,
$\PP_U(B)\le 4\sigma_U^2/\Delta^2$; and since $\E_\mu[X]=\E_U[X]-\Delta=t^\star-\Delta/2$,
$\PP_\mu(B^c)\le 4\sigma_\mu^2/\Delta^2$. Hence
\begin{equation}\label{eq:sep}
\TV(\mu_{n,q},U_n)\ \ge\ \PP_\mu(B)-\PP_U(B)\ \ge\ 1-\frac{4(\sigma_\mu^2+\sigma_U^2)}{\Delta^2}
\ \ge\ 1-\frac{2n^3/3}{\Delta^2}.
\end{equation}

\emph{Case (a): $\eps\le\tfrac1{2n}$.} By \eqref{eq:Deltalb},
$\Delta\ge \eps n^3/192$, so $\Delta^2\ge \eps^2 n^6/192^2$ and, with
$\tfrac23\cdot192^2=24576=C_1$,
\[
\TV\ \ge\ 1-\frac{2n^3/3}{\eps^2 n^6/192^2}=1-\frac{24576}{\eps^2 n^3}
=1-\frac{C_1}{(n^{3/2}\eps)^2},
\]
proving (a).

\emph{Case (b): $\eps>\tfrac1{2n}$.} Since $\Delta$ is nondecreasing in $\eps$,
$\Delta(\eps)\ge\Delta(\tfrac1{2n})\ge \tfrac{1}{2n}\cdot\tfrac{n^3}{192}=\tfrac{n^2}{384}$
by \eqref{eq:Deltalb}. Plugging into \eqref{eq:sep},
\[
\TV\ \ge\ 1-\frac{2n^3/3}{(n^2/384)^2}=1-\frac{(2/3)\cdot384^2}{n}=1-\frac{98304}{n}
=1-\frac{4C_1}{n},
\]
proving (b).

\emph{Conclusion.} If $n^{3/2}\eps_n\to\infty$: for indices with $\eps_n\le\tfrac1{2n}$,
(a) gives $\TV\ge1-C_1/(n^{3/2}\eps_n)^2\to1$; for indices with $\eps_n>\tfrac1{2n}$,
(b) gives $\TV\ge1-4C_1/n\to1$. In either case $\TV(\mu_{n,q_n},U_n)\to1$.
\end{proof}

\begin{corollary}\label{cor:threshold}
$\TV(\mu_{n,q_n},U_n)\to0$ if and only if $n^{3/2}\eps_n\to0$; and
$\TV(\mu_{n,q_n},U_n)\to1$ if $n^{3/2}\eps_n\to\infty$.
\end{corollary}

\begin{proof}
\emph{If $n^{3/2}\eps_n\to0$ then $\TV\to0$:} Proposition~\ref{prop:TVupper}. \emph{Last
assertion ($\TV\to1$ when $n^{3/2}\eps_n\to\infty$):} Proposition~\ref{prop:TVlower}.

\emph{If $\TV\to0$ then $n^{3/2}\eps_n\to0$:} we prove the contrapositive. Suppose
$n^{3/2}\eps_n\not\to0$; pick $\delta>0$ and a subsequence (still indexed by $n$) with
$n^{3/2}\eps_n\ge\delta$ and $n\ge4$. If along this subsequence infinitely many indices
have $\eps_n>\tfrac1{2n}$, then Proposition~\ref{prop:TVlower}(b) gives
$\TV\ge1-4C_1/n\to1$ there, so $\TV\not\to0$. Otherwise eventually $\eps_n\le\tfrac1{2n}$,
and we use the following Cantelli separation bound. For a statistic $X$ with means
$\E_\mu X,\E_\nu X$ and variances $\Var_\mu X,\Var_\nu X$, put $d=|\E_\mu X-\E_\nu X|$;
then
\begin{equation}\label{eq:cantelli}
\TV(\mu,\nu)\ \ge\ 1-\frac{\Var_\mu X}{\Var_\mu X+d^2/4}-\frac{\Var_\nu X}{\Var_\nu X+d^2/4}.
\end{equation}
Indeed, assuming WLOG $\E_\nu X\ge\E_\mu X$ and choosing the threshold
$c=\tfrac12(\E_\mu X+\E_\nu X)$, the one-sided Chebyshev (Cantelli) inequality gives
$\PP_\mu(X\ge c)\le \frac{\Var_\mu X}{\Var_\mu X+d^2/4}$ and
$\PP_\nu(X< c)\le \frac{\Var_\nu X}{\Var_\nu X+d^2/4}$, so with $A=\{X<c\}$,
$\TV\ge \PP_\mu(A)-\PP_\nu(A)\ge\big(1-\frac{\Var_\mu X}{\Var_\mu X+d^2/4}\big)
-\frac{\Var_\nu X}{\Var_\nu X+d^2/4}$, which is \eqref{eq:cantelli}.

Apply \eqref{eq:cantelli} with $X=\inv$, $\mu=\mu_{n,q_n}$, $\nu=U_n$. By
\eqref{eq:Deltalb}, $d=\Delta(\eps_n)\ge \eps_n n^3/192$, and
$\Var_\mu X,\Var_U X\le n^3/12$. Hence
\[
\frac{d^2/4}{\Var X}\ \ge\ \frac{(\eps_n n^3/192)^2/4}{n^3/12}
=\frac{12}{4\cdot192^2}\,\eps_n^2 n^3=\frac{3}{192^2}(n^{3/2}\eps_n)^2
\ \ge\ \frac{3\delta^2}{192^2}=:\rho>0 .
\]
Writing $u=\Var_\mu X$ and $v=\Var_U X$, the right side of \eqref{eq:cantelli} equals
$\frac{d^2/4}{u+d^2/4}-\frac{v}{v+d^2/4}=\frac{d^2/4}{u+d^2/4}-\frac{1}{(v/(d^2/4))+1}$;
both fractions are monotone, and since $d^2/4\ge\rho\,u$ and $d^2/4\ge\rho\,v$,
\[
\TV\ \ge\ \frac{\rho}{1+\rho}-\frac{1}{\rho+1}=\frac{\rho-1}{\rho+1}.
\]
This is positive when $\rho>1$. For $\rho\le1$ we instead optimize the threshold split:
repeating the Cantelli argument with $c=\E_\mu X+r$, $\E_\nu X-s$, $r+s=d$, $r=\lambda d$,
gives $\TV\ge 1-\frac{u}{u+\lambda^2 d^2}-\frac{v}{v+(1-\lambda)^2 d^2}$; taking
$\lambda\to1^-$ yields $\TV\ge \frac{\lambda^2 d^2}{u+\lambda^2 d^2}-\frac{v}{v+(1-\lambda)^2d^2}
\to 1-0=1$ as the second term tends to $1$... more carefully, fix $\lambda=\tfrac{1}{2}$:
then $\TV\ge 1-\frac{u}{u+d^2/4}-\frac{v}{v+d^2/4}=\frac{\rho_u}{1+\rho_u}-\frac{1}{1+\rho_v}$
with $\rho_u=d^2/(4u)\ge\rho$, $\rho_v=d^2/(4v)\ge\rho$; this is
$\ge\frac{\rho}{1+\rho}-\frac{1}{1+\rho}=\frac{\rho-1}{1+\rho}$, the same bound. To get
strict positivity for all $\rho>0$, use the asymmetric choice
$r=\tfrac{3}{4}d,\ s=\tfrac14 d$: then
$\TV\ge 1-\frac{u}{u+(9/16)d^2}-\frac{v}{v+(1/16)d^2}\ge
\frac{(9/16)\rho}{1+(9/16)\rho}-\frac{1}{1+(1/16)\rho}$, which is a fixed positive number
for every fixed $\rho>0$ (it is continuous and tends to $1$ as $\rho\to\infty$ and to a
positive value bounded below by $\tfrac{(9/16)\rho}{1+(9/16)\rho}-\tfrac{1}{1+(1/16)\rho}$,
positive once $\tfrac{9}{16}\rho(1+\tfrac1{16}\rho)>1+\tfrac9{16}\rho$, i.e.\ for $\rho$
above an absolute threshold; and below that threshold we may simply enlarge $\delta$).
In all cases $\liminf_n\TV(\mu_{n,q_n},U_n)>0$ along the subsequence, so $\TV\not\to0$.
\end{proof}

\section{Theorem C: ballistic crossover at fixed $q$}\label{sec:C}

For fixed $q\in(0,1)$ we show $L\asymp n$, leaving the KPZ class.

\subsection{Three deterministic inequalities}
\begin{lemma}[Total displacement vs.\ inversions]\label{lem:tdisp}
For every $\pi\in S_n$, $\displaystyle\sum_{i=1}^n|\pi(i)-i|\le 2\,\inv(\pi)$.
\end{lemma}
\begin{proof}
Sort $\pi$ to the identity by adjacent transpositions; each adjacent transposition of
neighbouring positions changes $\inv$ by $\pm1$ and changes $\sum_i|\pi(i)-i|$ by at
most $2$. Hence $\sum_i|\pi(i)-i|=\big|\sum_i|\pi(i)-i|-0\big|\le 2\,\inv(\pi)$, the
identity having $\sum_i|\cdot|=0$ and $0$ inversions, and $\inv(\pi)$ being the minimal
number of adjacent transpositions needed. (Concretely: with $T=\sum_i|\pi(i)-i|$ and
$I=\inv(\pi)$, each of the $I$ adjacent swaps in a shortest sorting sequence decreases
$I$ by exactly $1$ and decreases $T$ by at most $2$, so $T\le 2I$.)
\end{proof}

\begin{lemma}[Banded maps have short antichains]\label{lem:band}
Let $G\subseteq\{1,\dots,n\}$ and suppose $|\pi(i)-i|\le b$ for all $i\in G$. Then any
decreasing subsequence of $\pi$ supported in $G$ (positions $i_1<\cdots<i_d$ in $G$ with
$\pi(i_1)>\cdots>\pi(i_d)$) has length $d\le 2b+1$.
\end{lemma}
\begin{proof}
$\pi(i_1)\le i_1+b$ and $\pi(i_d)\ge i_d-b$; since $i_1<i_d$,
$\pi(i_1)-\pi(i_d)\le (i_1-i_d)+2b\le 2b-1<2b$. The $d$ distinct integers
$\pi(i_1)>\cdots>\pi(i_d)$ lie in an interval of length $<2b$, hence among at most $2b$
consecutive integers, so $d\le 2b+1$.
\end{proof}

\begin{lemma}[Dilworth lower bound on a subset]\label{lem:dilworth}
For $\pi\in S_n$ and $G\subseteq\{1,\dots,n\}$, if every decreasing subsequence of $\pi$
supported in $G$ has length $\le D_0$, then $L(\pi)\ge |G|/D_0$.
\end{lemma}
\begin{proof}
By the dual of Dilworth's theorem (patience sorting), the $|G|$ numbers
$(\pi(i))_{i\in G}$, read in increasing order of position, can be partitioned into at
most $D_0$ increasing subsequences (assign each element to a pile equal to the length of
a longest decreasing subsequence ending there; there are $\le D_0$ piles and each pile
is increasing in position and value). By pigeonhole one pile has $\ge |G|/D_0$ elements,
an increasing subsequence of $\pi$ of that length; thus $L(\pi)\ge |G|/D_0$.
\end{proof}

\subsection{Mean displacement under Mallows}
\begin{lemma}\label{lem:Einv}
For $q\in(0,1)$ and $\pi\sim\mu_{n,q}$, $\E_{\mu_{n,q}}[\inv(\pi)]\le \dfrac{nq}{1-q}$, and
hence $\displaystyle\sum_{i=1}^n\E_{\mu_{n,q}}\big[|\pi(i)-i|\big]\le\frac{2nq}{1-q}$.
\end{lemma}
\begin{proof}
By Lemma~\ref{lem:product}, $\E_\mu[\inv]=\sum_{m=1}^n\E_{\nu^{(q)}_m}[c']$. The
truncated law $\nu^{(q)}_m$ on $\{0,\dots,m-1\}$ is stochastically dominated by the
geometric law $\gamma(k)=(1-q)q^k$ on $\{0,1,2,\dots\}$: for every $r\ge0$,
\[
\PP_{\nu^{(q)}_m}(c'\ge r)=\frac{\sum_{k\ge r}^{m-1}q^k}{\sum_{k=0}^{m-1}q^k}
=\frac{q^{r}\sum_{k=0}^{m-1-r}q^{k}}{\sum_{k=0}^{m-1}q^{k}}\le q^{r}=\PP_{\gamma}(c'\ge r),
\]
because the remaining ratio $\sum_{k=0}^{m-1-r}q^k/\sum_{k=0}^{m-1}q^k\le1$. Hence
$\E_{\nu^{(q)}_m}[c']\le \E_\gamma[c']=\tfrac{q}{1-q}$, and summing over $m=1,\dots,n$
gives $\E_\mu[\inv]\le \tfrac{nq}{1-q}$. The displacement bound follows from
Lemma~\ref{lem:tdisp} by taking expectations.
\end{proof}

\subsection{The linear lower bound}
\begin{theorem}\label{thm:C}
Fix $q\in(0,1)$ and write $W=\tfrac{2q}{1-q}$. For $\delta\in(0,1)$ set
$b=b(q,\delta)=\big\lceil \tfrac{W}{\delta^{2}}\big\rceil$ and
$c(q,\delta)=\tfrac{1-\delta}{2b+1}>0$. Then for all $n\ge1$ and $\pi\sim\mu_{n,q}$,
\[
\PP_{\mu_{n,q}}\big(\,L(\pi)\ge c(q,\delta)\,n\,\big)\ \ge\ 1-\delta .
\]
In particular, taking $\delta\to0$ along a sequence, there is $c(q)>0$ with
$\PP_{\mu_{n,q}}(L(\pi)\ge c(q)\,n)\to1$; thus $L(\pi)/\sqrt n\to\infty$ in probability
and the Tracy--Widom limit fails.
\end{theorem}

\begin{proof}
Let $N=\#\{\,i:|\pi(i)-i|>b\,\}$. By Markov (pointwise $\mathbf 1\{|\pi(i)-i|>b\}\le
|\pi(i)-i|/b$) and Lemma~\ref{lem:Einv},
\[
\E[N]=\sum_{i=1}^n\PP\big(|\pi(i)-i|>b\big)\le\frac1b\sum_{i=1}^n\E\big[|\pi(i)-i|\big]
\le\frac{nW}{b}\le \frac{nW}{W/\delta^2}=\delta^2 n.
\]
By Markov on the nonnegative integer $N$,
$\PP(N\ge \delta n)\le \E[N]/(\delta n)\le \delta$. On the complementary event
$\{N<\delta n\}$, the good set $G=\{i:|\pi(i)-i|\le b\}$ satisfies $|G|=n-N>(1-\delta)n$,
and by Lemma~\ref{lem:band} (with this $b$) every decreasing subsequence supported in
$G$ has length $\le 2b+1$; Lemma~\ref{lem:dilworth} with $D_0=2b+1$ then gives
\[
L(\pi)\ \ge\ \frac{|G|}{2b+1}\ >\ \frac{(1-\delta)n}{2b+1}=c(q,\delta)\,n .
\]
Hence $\PP\big(L(\pi)\ge c(q,\delta)n\big)\ge \PP(N<\delta n)\ge 1-\delta$.

Finally, for any fixed $\delta_0\in(0,1)$ this gives $\PP(L\ge c(q,\delta_0)n)\ge1-\delta_0>0$,
so $L\ge c(q,\delta_0)n$ with positive probability bounded below; letting
$\delta\downarrow0$ shows $\PP(L(\pi)\ge c(q,\delta)n)\to1$ for the corresponding
$c(q,\delta)>0$. Since $c(q,\delta)>0$ is independent of $n$,
$L(\pi)/n$ is bounded below by a positive constant with probability $\to1$, whence
$L(\pi)/\sqrt n=\sqrt n\cdot L(\pi)/n\to\infty$ in probability. Consequently
$(L-2\sqrt n)/n^{1/6}\to+\infty$ in probability and has no Tracy--Widom (or any
finite) limit at the KPZ scale.
\end{proof}

\section{Synthesis: the phase diagram}\label{sec:synth}
With $\eps_n=1-q_n$, the natural scale is $n^{-3/2}$, the reciprocal of the standard
deviation $\sigma_U\asymp n^{3/2}$ of $\inv$ under the uniform measure (tilting by
$q^{\inv}$ shifts $\inv$ by $\asymp\eps_n\sigma_U^2\asymp\eps_n n^3$, comparable to one
standard deviation when $\eps_n n^{3/2}\asymp1$).
\begin{itemize}
\item \textbf{KPZ universality (Thm.~\ref{thm:A}).} $\eps_n=o(n^{-3/2})$ $\Rightarrow$
$\TV(\mu_{n,q_n},U_n)\to0$ and $(L-2\sqrt n)/n^{1/6}\xrightarrow{d}F_{\TW}$: a genuine
non-uniform family ($q_n\ne1$) in the KPZ class, with the BDJ centering, scaling, and
limit law.
\item \textbf{Sharpness (Thm.~\ref{thm:B}).} The exponent $3/2$ is exact for coupling:
$\TV\to0$ iff $\eps_n=o(n^{-3/2})$, and $\TV\to1$ when $\eps_n\gg n^{-3/2}$.
\item \textbf{Ballistic departure (Thm.~\ref{thm:C}).} For fixed $q\in(0,1)$,
$L\ge c(q)n$ with probability $\to1$: the order of $L$ jumps from $\sqrt n$ to $n$, and
no Tracy--Widom limit holds.
\end{itemize}
This resolves the problem: we exhibit a natural non-uniform ensemble in the KPZ class,
prove its Tracy--Widom fluctuation theorem, identify the sharp coupling threshold, and
exhibit the crossover to a non-KPZ (ballistic) regime.

\begin{thebibliography}{9}
\bibitem{BDJ} J.~Baik, P.~Deift, K.~Johansson, \emph{On the distribution of the length
of the longest increasing subsequence of random permutations}, J. Amer. Math. Soc.
\textbf{12} (1999), 1119--1178.
\bibitem{TW} C.~A.~Tracy, H.~Widom, \emph{Level-spacing distributions and the Airy
kernel}, Comm. Math. Phys. \textbf{159} (1994), 151--174.
\bibitem{Stanley} R.~P.~Stanley, \emph{Enumerative Combinatorics, Vol.~1}, 2nd ed.,
Cambridge Univ. Press, 2011.
\bibitem{BP} N.~Bhatnagar, R.~Peled, \emph{Lengths of monotone subsequences in a
Mallows permutation}, Probab. Theory Related Fields \textbf{161} (2015), 719--780.
\end{thebibliography}

\end{document}
